\chapter{Conclusion}
\label{chap:conclusion}

\section{Research questions revisited}
Answer each research question directly using findings already established in the
thesis. The conclusion should not introduce a new dataset, method or major argument.

\section{Contributions}
Return to the claims in \Cref{tab:contribution-map} and state the demonstrated
contribution, its evidence and its boundary. Remove any contribution that the completed
study does not actually substantiate.

\section{Recommendations and future research}
Prioritise a small number of feasible recommendations. Future research should arise
from the study's limitations or new questions exposed by the findings, not from a
generic statement that more research is needed.

\section{Closing statement}
End with a concise statement of what the reader should remember. This final paragraph
is an opportunity to reconnect the technical contribution to the larger problem
introduced in \Cref{chap:introduction}.

